Private health insurance companies in various countries are allowed to vary their services and pricing options based off various information about the customer and the risks involved. The so-called ``risk factors'' involved when profiling include age, gender, race, sexuality, smoking habits, income, body mass index, medical history and more~\cite{odihi20}. Certain forms of profiling have already been banned in parts of the world\footnote{For example genetic profiling is currently forbidden in the US~\cite{genetic-info-nondiscrimination}}, while others are still common. As the high price of healthcare services remains a big issue across the world\footnote{For instance, nearly $2$ million people declare bankruptcy in the US because of medical bills every year~\cite{medical-bills-us-bankrupcy}}, the implications of statistical profiling are an important subject of discussion. In this article, we aim to analyse this practice and its various implications through both an ethical and a policy-based framework.
